\documentclass[10pt]{article}

\usepackage[utf8]{inputenc}

\usepackage[T1]{fontenc}

\usepackage{amsmath}

\usepackage{amsfonts}

\usepackage{amssymb}

\usepackage[version=4]{mhchem}

\usepackage{stmaryrd}

\usepackage{graphicx}

\usepackage[export]{adjustbox}

\graphicspath{ {./images/} }



\begin{document}

ровки $W_{2}^{1}(\Omega)$ получена (см. [2]) модификация Ф. н. вида



\[

\int_{\Omega} f^{2} d \Omega \leqslant C\left\{\int_{\Omega} \Sigma_{i=1}^{n}\left(\frac{\partial f}{\partial x_{i}}\right)^{2} d \Omega+\left(\int_{\Gamma} f d \Gamma\right)^{2}\right\} .

\]



Имеются обобщения (см. [3] - [5]) Ф. н. на весовые классы (см. Весовое пространство, Вложения теоремь). Пусть $\Gamma \subset C^{(l)}$, числа $r, p, \alpha$ действительные, причем $r$-натуральное, $1 \leqslant p<\infty$. Говорят, что $f \in W_{p, \alpha}^{r}(\Omega)$, если конечна норма



где



\[

\|f\|_{W_{p, \alpha}^{r}(\Omega)}=\|f\|_{L_{p}(\Omega)}+\|f\|_{\omega_{p, \alpha}^{r}(\Omega)}

\]



\[

\begin{gathered}

\|f\|_{L_{p}(\Omega)}=\left(\int_{\Omega}|f|^{p} d \Omega\right)^{1 / p}, \\

\|f\|_{\omega p, \alpha}^{r}(\Omega)=\sum_{|k|=r}\left\|\rho^{\alpha} f^{(k)}\right\|_{L_{p}(\Omega)}, \\

f^{(k)}=\frac{\partial^{|k|_{f}}}{\partial x_{1}^{k_{1}} \ldots \partial x_{n}^{k_{n}}}, \quad|k|=\sum_{i=1}^{n} k_{i},

\end{gathered}

\]



$\rho \equiv \rho(x)$-функция, эквивалентвая расстоянию от $x \in \Omega$ до $\Gamma$.



Пусть число $s_{0}$-натуральное и



\[

r-\alpha-1 / p \leqslant s_{0}<r-\alpha+1-1 / p .

\]



Тогда, если $\Gamma \subset C\left(s_{0}+1\right),-1 / p<\alpha<r-1 / p, r / 2 \leqslant s_{0}$, то для $f \in W_{p, \alpha}^{r}(\Omega)$ справедливо неравенство



\[

\|f\|_{L_{p}(\Omega)} \leqslant

\]



$\leqslant C\left\{\sum_{l+s<r / 2}\left\|\left(\left.\frac{\partial^{s} f}{\partial n^{s}}\right|_{\Gamma} ^{p}\right)^{(l)}\right\|_{L_{p}(\Gamma)}+\|f\|_{\omega_{p, \alpha}^{r}(\Omega)}\right\}$, где $\left.\frac{\partial s f}{\partial n^{s}}\right|_{\Gamma}$-производная порядка $s$ по внутренней нормали к $\Gamma$ в точках $\Gamma$.



Также получается и неравенство типа неравенства (2), к-рое в простейшем случае имеет вид



где



\[

\|f\|_{L_{p}(\Omega)}^{p} \leqslant C\left(\|f\|_{\omega_{p, \alpha}^{1}(\Omega)}^{p}+\left|\int_{\Gamma} u \tau d \Gamma\right|^{p}\right)

\]



\[

\begin{gathered}

p>1, \gamma>1,-1 / p<\alpha<1-1 / p-1 / \gamma, \\

\tau \in L \gamma(\Gamma), \int_{\Gamma} \tau d \Gamma \neq 0 .

\end{gathered}

\]



Всюду постоянная $C$ не зависит от $f$.



Неравенство названо по имени К. Фридрихса, к-рый получил его при $n=2, f \in C^{(2)}(\bar{\Omega})$ (см. [1]).



Jum.: [1] F r i e d r i c h s K., «Math. Ann.»), 1927, Bd 98,

$566-75 ;[2]$ C o 0 o л e в C. JI., Некоторые применения Функционального анализа в математической физике, Новосиб. 1962 ; [3] Н и к о л в с к и й С. М., Л и з о р к и н П. И., "докл. AH CCCP", 1964 , т. 159 , № 3, c. $512-15 ;$ [4] Н'и к о л bс к и й С. М., Приближение функий многих перемениых и теоремы вложения, 2 изд., М., 1977; [5] К а л и н и ч е н к о II. Ф.,,

«Матем. сб.», 1964 , т. 64 , № 3, с. $436-57 ;[6]$ К у р а н т Р., Г и л ь б е р т Д., Методы математической физики, пер. с нем., 2 изд., т. 2, М. - Л., 1951, [7] N i r e n b e r g L., "Ann. Scuola "Meddel. fràn Lunds univers. matem. semin.", 1955 , Bd $13, \mathrm{p}$. 84. Д. Ф. Калиниченко, Н. В. Мирошин. ФРОБЕНИУСА АВТОМОРФИЗМ - Элемент группы Галуа специального вида, играющий фундаментальную роль в теории полей классов. Пусть $L$-алгебраич. расширение конечного поля $K$. Тогда Ф. а. наз. автоморфизм $\varphi_{L / K}$, определяемый формулой $\varphi_{L / K}(a)=a^{q}$ для всех $a \in L$, где $q=\# K$ (мощность $K$ ). Если $L / K-$ конетное расширение, то $\varphi_{L / K}$ порождает группу Галуа $G(L / K)$. Для бесконетного расширения $L / K$ автоморфизм $\varphi_{L / K}$ является топологич. образующей группы $G(L / K)$. Если $L \supset E \supset K$ и $[E: K]<\infty$, то $\varphi_{L / E}=\varphi_{L / K}[E: K]$.



Пусть $k$-локальное поле с конечным полем вычетов $\bar{k}$, а $K$-неразветвленное расширение поля $k$. Тогда Ф. а. $\varphi_{\bar{K} / \bar{k}}$ расширевий полей вычетов однозначно продолжается до автоморфизма $\varphi_{K / k} \in G(K / k)$, наз. Ф. а. не р азветвленн о го р асти и ения $K / k$. Пусть $\# \bar{k}=q, \mathcal{O}_{K}$-кольцо целых элементов поля $K$ и $\mathfrak{p}$-максимальный идеал в $G_{K}$. Тогда Ф. а. $\varphi_{K / k}$ однозначно огределяется условием $\varphi_{K / k}(a) \equiv a^{q} \bmod \mathfrak{p}$ для любого $a \in \hat{O}_{K}$. Если $K / k-$ произвольное расширение Галуа локальных полей, то Ф. а. расширения $K / k$ иногда называют любой автоморфизм $\varphi \in G(K / k)$, индуцирующий на максимальном неразветвленном подрасширении поля $K$ Ф. а. в указанном выше смысле. Пусть $K / k$-расширение Галуа глобальных полей, ㄱ-простой идеал поля $k$ и $\mathfrak{S}$-нек-рый простой идеал поля $K$, лежащий над $\mathfrak{p}$. И пусть $\mathfrak{P}$ не разветвлен в расширении $K / k$ и $\varphi_{\mathfrak{k}} \in G\left(K_{\mathfrak{B}} / k_{\mathfrak{p}}\right)-\Phi$. а. неразветвленного расширения локальных полей $K_{\text {æ }} / k_{\mathfrak{p}}$. Отождествляя группу Галуа $G\left(K_{\mathfrak{\Re}} / k_{\mathfrak{p}}\right)$ с подгруппой разложения идеала ऐ в $G(K / k)$, можно рассматривать $\varphi_{3}$ как элемент группы $G(K / k)$. Этот элемент наз. Ф. а., соответствующим простому идеалу $\mathfrak{P}$. Если $K / k-\kappa 0-$ нечное распирение, то согласно те ореме Чебот а р е в а о пл отн ости для любого автоморфизма $\sigma \in G(K / k)$ существует бесконечное число простых идеалов $\mathfrak{B}$, не разветвленных в $K / k$ таких, что $\sigma=\varphi_{\mathfrak{R}}$. Для



\includegraphics[max width=\textwidth, center]{2023_07_09_c0b6a459caa595364489g-1}

$\mathfrak{p}$. В этом случае $\varphi_{\mathfrak{B}}$ обозначается через $(\mathfrak{p}, K / k)$ и наз. с и м в о лом А р т и на простого идеала $\mathfrak{p}$. Лит.: [1] В е й л ь А., Основы теории чисел, пер. с англ., M., $1972 . \quad$ Л. B. Кузз.миt. ФРОБЕНИУСА ТЕОРЕМА - теорема, описывающая все конечномерные ассоциативные действительные алгебры без делителей нуля, доказана Г. Фробениусом [1]. Ф. т. утверждает, что:



\begin{enumerate}

  \item Поле действительных чисел и поле комплексных чисел являются единственными конечномерными действительными ассоциативно-коммутативными алгебрами без делителей нуля.



  \item Тело кватернионов является единственной конөчномерной действительной ассоциативной, но не коммутативной алгеброй без делителей нуля.



\end{enumerate}



Существует также описание альтернативных конечномерных алгебр без делитөлей нуля:



\begin{enumerate}

  \setcounter{enumi}{2}

  \item Алгебра Кәли является единственной конечномерной действительной альтернативной, но не ассоциативной алгеброй без делителей нуля.

\end{enumerate}



Объединение этих трех утверждений наз. о б о бщенв о й те о рев ой участвующие в формулировке теоремы алгебры оказываются алгебрами с однозначным делением и с единицей Ф. т. не может быть обобщена на случай неальтернативных алгебр. Доказано, однако, что размерность любой конечномерной действительной алгебры без делителей нуля может принимать лишь значения, равные $1,2,4$ или 8.



Jum.. [1] F F o b e n i u s F., "J. reine und angew. Math.", 1877, Bd, 82, S. $230-315$, .

ллгебр, 2 пзд. M., 1973.



ФРОБЕНИУСА ТЕОРЕМА - теорема об условиях полной интегрируемости системы уравнений Пфаффа или (в геометрич. терминах) об условиях, при к-рых заданное ва дифференцируемом многообразии поле $n$-мерных касательных подпространств является касательным полем нек-рого слоенил. Несколько эквивалентных формулировок Ф. т. см. в статьях Инволютивное распределение, Коши задача; вариант с минимальными требовавиями дифференцируемости см. в [2]. Название Ф. т. связано с изложением этой теоремы в [1], но не соответствует ириводимым там сведениям об еө истории.





\end{document}